\documentclass{otscience}

\typeout{OT7C-STARTER:OTSCIENCE}

\title{\textbf{Starter Science Notes}}
\author{Author Name}
\date{\today}

\begin{document}

\maketitle
\tableofcontents
\newpage

\section{Concept}

\begin{definitionbox}[Quantity]
  A quantity is a measurable property of a system. It should be written with a
  value, a unit, and a short description of what was measured.
\end{definitionbox}

\begin{intuitionbox}[How to Read the Page]
  Start with the concept, write the controlling formula, then add one example
  that shows how the formula is used.
\end{intuitionbox}

\section{Formula}

\begin{formulabox}[Wave Speed]
  \[
    v=f\lambda
  \]
  where \(v\) is wave speed, \(f\) is frequency, and \(\lambda\) is wavelength.
\end{formulabox}

\begin{examplebox}[Worked Example]
  If \(f=\SI{5}{Hz}\) and \(\lambda=\SI{2}{m}\), then
  \[
    v=f\lambda=\SI{10}{m.s^{-1}}.
  \]
\end{examplebox}

\begin{warningbox}[Common Mistake]
  Do not mix units silently. Convert the inputs before substituting them.
\end{warningbox}

\begin{summarybox}[Takeaway]
  A good science note should let you recover the meaning, the formula, and one
  worked use without reopening the source lecture.
\end{summarybox}

\end{document}
